\section[Classical K-theory]{Classical $K$-theory}
Before expanding the tools from algebraic $K$-theory to the world of graphs, it is worth building up the classical theory to get a good idea of which ideas to use.

\subsection[The K-theory spectrum in many contexts]{The $K$-theory spectrum in many contexts}
\subsubsection[The K-theory of a ring]{The $K$-theory of a ring}
It is no secret that the dimension and the determinant are very useful tools in the study of vector spaces over a field. The low degree $K$-theory of a ring aims to capture this data for modules.

\begin{defn}
    Let $R$ be a ring. The \emph{$0$-th $K$-theory group} of $R$ is the abelian group 
    \[K_0(R) = \nicefrac{\Fab{\Iso{\Proj{R}}}}{\sim}\]
    where the equivalence relation is generated by $M \sim L + N$ whenever there is a short exact sequence 
    \[ 0 \to L \to M \to N \to 0 \]

    The \emph{$1$-st $K$-theory group} of $R$ is the abelian group
    \[K_1(R) = \nicefrac{\GL{R}}{\E{R}} \]
\end{defn}

Both groups may be obtained via alternate constructions which we will not detail. We only recall a few elementary properties of both to motivate their existence.

The $0$-th $K$-theory group serves as a universal space for the rank of a projective module. A \emph{generalized rank} on $\Proj{R}$ consists of an abelian group $(A, +, 0)$ with a morphism
\[ r: \Iso{\Proj{R}} \to A \]
such that $r(M) = r(L) + r(N)$ whenever there is a short exact sequence
\[ 0 \to L \to M \to N \to 0.\]
The way $K_0(R)$ is constructed allows it to be a universal generalized rank, meaning that the natural map $\tilde{r} : \Iso{\Proj{R}} \to K_0(R)$ is a generalized rank, such that for any other generalized rank $r : \Iso{\Proj{R}} \to A$, there is a unique group homomorphism $\hat{r} : K_0(R) \to A$ such that the following diagram is commutative:
% https://q.uiver.app/#q=WzAsMyxbMCwwLCJcXElzb3tcXFByb2p7Un19Il0sWzEsMCwiS18wKFxcUHJvantSfSkiXSxbMCwxLCJBIl0sWzAsMiwiciJdLFswLDEsIlxcdGlsZGV7cn0iLDJdLFsxLDIsIlxcaGF0e3J9Il1d
\[\begin{tikzcd}
	{\Iso{\Proj{R}}} & {K_0(\Proj{R})} \\
	A
	\arrow["{\tilde{r}}"', from=1-1, to=1-2]
	\arrow["r", from=1-1, to=2-1]
	\arrow["{\hat{r}}", from=1-2, to=2-1]
\end{tikzcd}\]

If $R$ is actually a field, there is an isomorphism $K_0(R) \cong \ZZ$, so that it happens to be the group completion of $\NN$, and the generalized rank is precisely given by the dimension of a vector space.

 Similarly, the group $K_1(R)$ comes with the quotient map $\tilde{\delta}: \GL{R} \to K_1(R)$ which is a universal generalized determinant.


\subsubsection{Waldhausen categories}
